% DRAFT. Every number comes from bench/results/*.md in this repository;
% see paper/README.md for the mapping. 
\documentclass[11pt]{article}
\usepackage[margin=1in]{geometry}
\usepackage[T1]{fontenc}
\usepackage{lmodern}
\usepackage{booktabs}
\usepackage{amsmath}
\usepackage{natbib}
\usepackage[hidelinks]{hyperref}
\usepackage{enumitem}
\usepackage{listings}
\setlist{nosep}
\emergencystretch=1.5em
\lstset{basicstyle=\ttfamily\small,columns=fullflexible,keepspaces=true,
  frame=single,framesep=4pt,xleftmargin=6pt,xrightmargin=6pt,
  showstringspaces=false}

\title{Myelin-AI: A Small Memory Layer for Language Agents\\
That Dates Its Facts, Learns From Use, and Forgets Reversibly}
\author{Scott Benson\\Apex Insights}
\date{DRAFT --- November 2026}

\begin{document}
\maketitle

\begin{abstract}
\mbox{Myelin-AI} is an open-source memory layer for language agents: a Python
library over one SQLite file, with no required dependencies and no
model of its own. It stores each message or fact with the time from
which it holds, retrieves with BM25 and optional CPU-only embeddings,
returns as much as fits a token budget, and keeps three kinds of state
that most retrieval stores do not: which memories answers actually
used, how memories relate to one another, and a single health grade
from which forgetting is decided. Forgetting removes a memory from
search and keeps it in the store, so it can be undone. With a
6{,}000-token context \mbox{Myelin-AI} scores 90.1\% on LongMemEval-S
(\texttt{gpt-5-mini}, mean of four runs; 88.6\% with \texttt{gpt-4o})
and 92.6\% on LoCoMo, within the range that published systems report,
with median retrieval latency of 30\,ms at 50{,}000 memories. An
optional embedder fine-tuned on synthetic conversations puts all the
evidence in the top five results for 66.7\% of LoCoMo questions and
72.4\% of LongMemEval-S questions, against 46.4\% and 65.1\% for the
default, and embeds about 150 times more slowly on a CPU. We
describe the design, say which parts moved those scores and which did
not, and report where it fails: on long histories whose old facts are
asked for once and late, forgetting by age hides evidence that use has
had no chance to protect. A companion paper gives the longitudinal
evaluation under which the feedback and forgetting mechanisms were
measured. Code, harnesses and result files are public under the MIT
licence.
\end{abstract}

\section{Introduction}

An agent that works with someone for months needs to remember what was
said, know which version of a changing fact is current, and avoid
carrying everything forever. The usual building block is a vector
store with a similarity search in front of it. That handles the first
need and none of the others: every item is equally current, nothing is
learned from which results were useful, and the store only grows.

\mbox{Myelin-AI} is a library that takes those needs as its design brief while
staying small enough to embed. (The name borrows from the myelin sheath
that speeds signals along frequently used nerve fibres; we write
\mbox{Myelin-AI} throughout to keep the software distinct from the
biology.) It has four commitments.

\begin{itemize}
\item \textbf{Facts carry time.} Every memory has a time from which it
  is valid and, once superseded, a time from which it is not. Search
  can be asked what was true at a given moment.
\item \textbf{Use is recorded.} After an answer, the caller reports
  which retrieved memories it relied on. That changes ranking for later
  questions and protects those memories from being forgotten.
\item \textbf{Forgetting is one decision, and reversible.} A single
  health grade combines age, confidence, contradictions, use and the
  use of related memories. Memories below a threshold are quarantined:
  out of search, still in the store, restored if the grade recovers.
\item \textbf{No model inside.} The core imports only the standard
  library. Embeddings, fact extraction and answering are optional and
  take whatever model the caller supplies.
\end{itemize}

This paper describes the system as released in version 1.5.0, with
the optional embedder added in 1.6.0, reports
its measured behaviour, and is explicit about what the measurements do
not support. It makes no claim of rank among memory systems; published
scores use different answering models, judges and prompts, and our own
identical runs differ by more than the gaps between systems near the
top.

\section{Related work}

MemGPT \citep{packer2023memgpt} gives an agent tools to page its own
memory in and out of context. Generative Agents
\citep{park2023generative} rank memories by recency, importance and
relevance. MemoryBank \citep{zhong2024memorybank} applies a forgetting
curve after \citet{ebbinghaus1885}. Mem0 \citep{chhikara2025mem0}
extracts and reconciles facts with a model at write time, Zep
\citep{rasmussen2025zep} builds a temporal knowledge graph in which
contradicted facts are invalidated, A-MEM \citep{xu2025amem} links
notes to one another, and Hindsight \citep{hindsight2025} and Mastra's
observational memory \citep{mastra2026om} organise what is retained
around model-written observations.

\mbox{Myelin-AI} differs mainly in what it leaves out. It makes no model call on
the write path unless asked, stores what was said and not a rewrite of
it, and keeps its state in ordinary SQL tables. Its health grade is in
the tradition of base-level and spreading activation in ACT-R
\citep{anderson2004actr}, much simplified, with one rule added for
agents: no amount of use may outweigh a contradiction.

\section{Design}

\subsection{Data model}

A memory is a row: text, a user scope, \texttt{t\_valid} (when the
fact began to hold), \texttt{t\_invalid} (when it stopped, if it has),
a confidence and an importance in $[0,1]$, free-form metadata, a
health state, and counters for use. A conversation is stored one
memory per message, with its session and position in the metadata.
Nothing is rewritten at write time; the stored text is what was said.

Two times are kept apart: when a fact holds in the world and when the
store learned it. \texttt{search(query, as\_of=t)} returns what was
valid at $t$, so a superseded address is still findable for the period
in which it was current, and is not returned for a question about now.

\subsection{Retrieval}

Lexical search is BM25 \citep{robertson2009bm25} over a SQLite
full-text index with stemming. If an embedder is supplied, cosine
similarity over its vectors is a second ranking, and the two are merged
by reciprocal-rank fusion \citep{cormack2009rrf}. The bundled embedder
is a static model that runs on a CPU; any object with an
\texttt{encode} method can replace it, and if it raises, reads and
writes continue lexically. An optional fine-tuned transformer embedder
is described in Section~\ref{sec:embedder}.

One fusion did not serve both the top of the list and its depth.
Weighting the vector ranking double put the most evidence inside a
6{,}000-token context on LongMemEval-S, and was worse than lexical
search alone in the top ten on LoCoMo's short messages (39.6\% of
questions with all evidence in the top ten, against 54.9\%). \mbox{Myelin-AI}
orders the first 20 results by an equal-weight fusion with a small
rank constant and the rest by the vector-heavy one. That gave 51.8\%
in the top ten on LoCoMo and left LongMemEval-S unchanged.

\subsection{Context under a budget}

\texttt{recall(query, token\_budget)} returns as many of the best
results as fit. A memory longer than a limit (100 tokens by default)
counts as an excerpt of the passages that share words with the query;
the three best results stay whole. On LongMemEval-S, lowering that
limit from 300 to 100 tokens raised the share of questions with every
evidence turn inside a 6{,}000-token context from 92.3\% to 97.0\%:
evidence is mostly in short user messages, and long assistant replies
were spending the budget. \texttt{format\_context} lays results out
oldest first under date headings, and \texttt{build\_answer\_prompt}
adds rules for the cases answering models get wrong: values that
changed, totals and orderings, relative dates, and questions the
history does not cover.

\subsection{Facts at write time}

\texttt{add\_conversation} can take a model. It then makes one call
per conversation to extract short dated facts (``Ana adopted a dog on
2023-05-19'') and indexes them with the messages they came from. The
message remains what is stored and shown; the fact is a second way to
find it. This is the only model call \mbox{Myelin-AI} can make on the write
path, and it is off by default.

\subsection{Feedback from use}

After answering, the caller passes \texttt{cite\_used} the results it
was shown and the ones it used. The numbered context and prompt let
the answering model supply that list itself, as a final line naming
entries. Memories that were used are counted as useful. Memories that
were shown and not used are counted too; this lowers their rank the
next time they match and has no effect on their health. An earlier
interface, \texttt{cite(id, answer\_correct)}, records whether an
answer that relied on a memory turned out right or wrong.

\subsection{One health grade}
\label{sec:health}

Each memory has a health $h=\sigma(r)$ with
\begin{align*}
r ={}& b - a(\text{age})
       - w_c\,(\text{contradictions})
       - w_q\,(1-\text{confidence})
       - w_i\,(1-\text{importance}) \\
     &+ \min\bigl(L,\;
        w_u \ln(1+u) + w_g \ln(1+s) + w_t\,\rho\bigr)
       - w_x \ln(1+x).
\end{align*}
Here $a(\cdot)$ grows once a memory is more than thirty days old, $u$
counts the times the memory was cited as used, $\rho$ decays with the
time since its last use, $x$ counts cites in answers that were wrong,
and $s$ is support from related memories. Support flows from each used
memory to its neighbours: messages within three positions in the same
conversation, weighted by $1/(1+\text{distance})$, and its eight
nearest memories by embedding with cosine similarity of at least 0.5,
weighted by that similarity, both scaled by $\ln(1+u)$ of the source.

The cap $L$ is set just below $w_c$, so one contradiction outweighs
any amount of use, direct or inherited. A frequently retrieved false
statement is the memory an agent can least afford to protect. The
inequality is checked when the library is imported and by a test.

\subsection{Forgetting}

\texttt{prune} first marks values that a later fact supersedes as no
longer current, setting \texttt{t\_invalid} without deleting them. It
then computes health and quarantines memories below 0.4. A quarantined
memory is absent from search and present in the store; it returns if
its grade recovers and can be inspected at any time. Deletion happens
only on explicit request, after a waiting period. In releases before
1.5.0, \texttt{prune} deleted; that was changed because a forgetting
rule that cannot be audited or undone cannot be tuned safely.

\subsection{Other components}

The library also contains a salience gate that scores a memory's
importance at write time, and a consolidation routine that tests the
store against held-out questions about its own contents during idle
time. Both predate the work reported here. Neither is used in the
measurements below, and we make no claim about their effect.

\section{Interface}

The whole path from history to answer is a few calls:

\begin{lstlisting}[language=Python]
from myelin import (MemoryClient, build_answer_prompt, format_context,
                    parse_used)
from myelin.embedders import StaticEmbedder

client = MemoryClient(db_path="memory.db", embedder=StaticEmbedder())
client.add_conversation(turns, when=session_date)   # one memory per message

hits = client.recall(question, token_budget=6000)
prompt = build_answer_prompt(format_context(hits, numbered=True),
                             question, ask_used=True)
reply = your_model(prompt)                          # any provider

client.cite_used(hits, parse_used(reply))           # what the answer used
client.prune()                                      # supersede, then quarantine
\end{lstlisting}

An optional HTTP sidecar exposes the same operations (store, search,
recall, cite, prune and others) for agents not written in Python.

\section{Evaluation}
\label{sec:eval}

All runs use hybrid search with the bundled static embedder and a
6{,}000-token context. LongMemEval-S \citep{wu2025longmemeval} is
judged by \texttt{gpt-4o} with the benchmark's prompts. LoCoMo
\citep{maharana2024locomo} is judged by \texttt{gpt-4o-mini} with the
prompt that published scores use.

\begin{table}[t]
\centering
\small
\begin{tabular}{lrr}
\toprule
 & \texttt{gpt-5-mini} (mean of 4) & \texttt{gpt-4o} (1 run) \\
\midrule
Overall & 90.1 & 88.6 \\
Single-session, user & 96.4 & 97.1 \\
Single-session, assistant & 96.4 & 94.6 \\
Single-session, preference & 85.0 & 76.7 \\
Multi-session & 79.9 & 80.5 \\
Temporal reasoning & 92.7 & 89.5 \\
Knowledge update & 94.9 & 93.6 \\
\bottomrule
\end{tabular}
\caption{LongMemEval-S accuracy by question type (\%). The four
\texttt{gpt-5-mini} runs scored 90.2, 91.0, 89.4 and 89.8.}
\label{tab:lme}
\end{table}

\paragraph{LongMemEval-S.} Table~\ref{tab:lme} gives accuracy by
question type. Release 1.4.0 scored 60.4\% with \texttt{gpt-4o-mini}.
All evidence turns are inside the context for 96.4\% of questions.
With the evidence sessions handed to the model directly, bypassing
retrieval, \texttt{gpt-5-mini} scores 93.4\%, so about three points
remain between this system and a perfect retriever with the same
reader. Most of that is in multi-session questions that require
counting items across sessions.

\paragraph{LoCoMo.} With write-time facts, 92.6\% of 1{,}540 questions
are judged correct (single-hop 95.7, multi-hop 90.8, temporal 91.0,
open-domain 76.0; one run). Without facts the figure is 89.0\%. An
independent audit reports that 6.4\% of LoCoMo's answer key is wrong
and that this judge prompt is lenient \citep{penfield2026locomo}, which
puts the ceiling near 93.6\% and means scores at this level should be
read with caution.

\paragraph{What moved the scores.} In rough order of size on
LongMemEval-S: the answering prompt and its rules; the answering
model; ranked lexical search with a token budget in place of a fixed
top ten; query-focused excerpts; and hybrid search. Write-time facts
added about three points on LoCoMo and nothing measurable on
LongMemEval-S. Extra queries written by a model, larger budgets, a
model-based filter and a stricter counting rule did not help.

\paragraph{Noise.} Runs of an identical configuration differed by up
to 3.4 points, and about 40 of 500 answers change between two such
runs. Eight runs across the last two configurations average 90.7\%.
We did not adopt changes on the strength of a single run.

\paragraph{For orientation.} Published LongMemEval-S scores include
94.9 with \texttt{gpt-5-mini} and 84.2 with \texttt{gpt-4o}
\citep{mastra2026om} and 91.4 with a larger model
\citep{hindsight2025}; on LoCoMo, systems' authors report between 74.0
and 94.7. These are self-reported with differing models, judges and
prompts, and are not a like-for-like comparison.

\begin{table}[t]
\centering
\small
\begin{tabular}{rrrrr}
\toprule
 & \multicolumn{2}{c}{Lexical} & \multicolumn{2}{c}{Hybrid} \\
Memories & p50 & p95 & p50 & p95 \\
\midrule
5{,}000 & 3 & 13 & 8 & 17 \\
50{,}000 & 15 & 62 & 30 & 78 \\
200{,}000 & 45 & 190 & 83 & 367 \\
\bottomrule
\end{tabular}
\caption{\texttt{recall} latency in milliseconds for one user's store
on synthetic text, one process, one SQLite file. Machine-dependent.}
\label{tab:scale}
\end{table}

\paragraph{Cost and scale.} Table~\ref{tab:scale} gives retrieval
latency as a store grows. Writes run at about 17{,}000 to 21{,}000
memories per second lexically and 2{,}300 to 3{,}700 with embeddings.
Latency grows in proportion to the store, since hybrid search compares
the query with every stored vector; an approximate index would be the
next step beyond a few hundred thousand memories per user. Vectors take
about 5\,KB per memory. Storing a conversation makes no model call
unless fact extraction is requested.

\section{Behaviour over time}

Feedback, relationships and forgetting do not change the scores above,
because those benchmarks ask each question once of a freshly written
store. A companion paper \citep{myelin2026longitudinal} describes
LongHAUL (Long-horizon Assessment of Use and Loss), a protocol that replays dated conversations on a simulated clock, asks
each question some sessions after its evidence, and scores retrieval
without a model. We summarise its results for this system.

\begin{table}[t]
\centering
\small
\begin{tabular}{lrrr}
\toprule
 & All evidence & Needed & Store still \\
 & in top 5 & but hidden & searchable \\
\midrule
Store and search only & 46.3 & 0.0 & 100.0 \\
Forgetting by age & 48.9 & 8.6 & 66.2 \\
Feedback & 52.7 & 0.0 & 100.0 \\
Feedback and forgetting & 52.8 & 2.5 & 86.2 \\
\bottomrule
\end{tabular}
\caption{LongHAUL on LoCoMo: each question asked up to six
sessions after its evidence (\%, three schedules).}
\label{tab:time}
\end{table}

\paragraph{Feedback helps.} On LoCoMo, reporting which memories were
used raises the share of questions with all evidence in the top five
from 46.3\% to 52.7\% (Table~\ref{tab:time}), and from 44.1\% to
53.7\% when questions are asked again. The gain is the same when
\texttt{gpt-4o-mini} names what it used in place of the evidence
labels (52.7\%). It carries to LongMemEval-S (65.1\% to 69.4\%).

\paragraph{Forgetting by age alone is costly.} It removes a third of
the store from search and hides a memory that a later question needs
for 8.6\% of questions. No setting of decay speed and threshold
avoided this.

\paragraph{Graded together, the cost falls.} With use and related use
in the grade, needed-but-hidden falls to 2.5\% while 14\% of the store
is still pruned.

\paragraph{Where it fails.} On the 120 LongMemEval-S histories that
span a month or more, forgetting hides a needed memory for 11.2\% of
askings and feedback recovers little (9.8\%). Each of those histories
has one question, so its old evidence has never been retrieved when
decay reaches it. Use can protect only what has been used. Part of the
measured cost may be correct behaviour, since most affected questions
concern a value that changed and the labels count the superseded turn
as needed; we did not measure answers. Either way, an application
whose users ask about old, never-revisited details should run \mbox{Myelin-AI}
with quarantine disabled or with a slower decay, and this is the main
open problem in the design.

\section{A fine-tuned embedder}
\label{sec:embedder}

The static embedder averages token vectors. It is fast and it is the
weakest part of retrieval: on LoCoMo, fewer than half of the questions
have all their evidence in the top five. From version 1.6.0
\mbox{Myelin-AI} offers a second, optional embedder, a small
transformer fine-tuned for one task: given a question, find the turn or
note that answers it. This section reports what it gains and what it
costs.

\paragraph{Model and data.} The base is \texttt{bge-small-en-v1.5}
\citep{xiao2023bge}, a 33M-parameter English model with 384-dimensional
vectors. The training data is synthetic. A language model wrote 1{,}000
items from seeded briefs: 500 chats between two people, 260 histories
of one person talking to an assistant, and 240 sets of the notes a
working agent saves about a project. Together they hold about 89{,}000
turns and notes and 20{,}000 questions, each tagged with the text that
answers it. From 900 of the items we built 23{,}781 training examples,
a question with its answering text and three near misses from the same
conversation, chosen as the static embedder's closest wrong results.
Training was one pass with a multiple-negatives ranking loss
\citep{henderson2017mnrl} in sentence-transformers
\citep{reimers2019sbert}, and took under two hours on a four-core CPU
with no GPU. No text from LoCoMo or LongMemEval was used. On the
held-out tenth of the items the answering text was the top result for 59.3\% of
questions, against 43.8\% before training.

\begin{table}[t]
\centering
\small
\begin{tabular}{llrrr}
\toprule
 & Embedder & Top 5 & Top 10 & Top 20 \\
\midrule
LoCoMo & Static (default) & 46.4 & 55.8 & 63.9 \\
 & \texttt{bge-small}, as published & 56.5 & 66.1 & 74.4 \\
 & Fine-tuned & 66.7 & 75.7 & 82.0 \\
\midrule
LongMemEval-S & Static (default) & 65.1 & 77.9 & 86.6 \\
 & \texttt{bge-small}, as published$^\dagger$ & 62.7 & 80.1 & 89.2 \\
 & Fine-tuned & 72.4 & 85.4 & 92.1 \\
\bottomrule
\end{tabular}
\caption{Share of questions with all evidence in the top results (\%),
LongHAUL with nothing forgotten and no feedback. LoCoMo: ten
conversations, three schedules. LongMemEval-S: 500 questions, one
schedule. $^\dagger$On a 250-question sample drawn evenly across
question types, where the static embedder scores 64.9 and the
fine-tuned one 72.2 in the top five.}
\label{tab:embedder}
\end{table}

\paragraph{Lift.} Table~\ref{tab:embedder} gives retrieval with each
embedder. On LoCoMo the published model is worth ten points over the
static one at every depth and fine-tuning adds ten more. On
LongMemEval-S the published model is no better than the static one in
the top five; the fine-tuned model is 7.3 points better there and 5.5
to 7.5 better at ten and twenty.

\paragraph{The training data has to resemble the use.} A first round
trained on the two-person chats alone. It reached 66.0\% in the top
five on LoCoMo, whose conversations are of that kind, and 65.1\% on the
LongMemEval-S sample, level with the static embedder. Adding the
assistant histories and the notes raised the sample to 72.2\% and left
LoCoMo where it was (66.7\%). The second round bought nothing on the
data the first round already covered and seven points on the data it
did not.

\begin{table}[t]
\centering
\small
\begin{tabular}{lrr}
\toprule
 & Static & Fine-tuned \\
\midrule
Short messages embedded per second & 14{,}914 & 103 \\
Long messages embedded per second & 5{,}007 & 15 \\
One question embedded (ms, median) & 0.09 & 11.7 \\
Time to load (s) & 1.4 & 4.4 \\
Weights on disk (MB) & 125 & 128 \\
Vector dimensions & 512 & 384 \\
Runtime required & numpy & PyTorch \\
\bottomrule
\end{tabular}
\caption{Cost of each embedder on a four-core CPU with no GPU, in
batches of 32. Short messages average about 140 characters and long
ones 1{,}220. Machine-dependent.}
\label{tab:embedcost}
\end{table}

\paragraph{Time.} Table~\ref{tab:embedcost} gives the cost. The
fine-tuned model embeds about 150 times fewer short messages a second
than the static one, and over 300 times fewer long ones. Writing is
where this is felt: with the static embedder a store takes 2{,}300 to
3{,}700 memories a second, and with the fine-tuned one it is limited
to about a hundred short messages a second, so importing a history of
500 long messages takes about half a minute in place of well under a
second. Reading is affected less. Embedding the question adds about
12\,ms to a \texttt{recall} whose median is 8 to 83\,ms
(Table~\ref{tab:scale}), and the vectors are a quarter smaller. The model also brings PyTorch, about 800\,MB
installed, into a package that otherwise needs only numpy for hybrid
search. We did not measure a GPU.

\paragraph{Feedback matters less with a better embedder.} On LoCoMo,
feedback adds 6.3 points in the top five with the static embedder, 3.5
with the published model and 1.5 with the first-round fine-tune (it
was not measured with the final model). The two mechanisms recover
some of the same misses. Forgetting by age costs about the same
with either embedder (2.5 and 2.7 points in the top five on
LongMemEval-S).

\paragraph{What this does not show.} The accuracy figures in
Section~\ref{sec:eval} were measured with the static embedder; answers
with the fine-tuned one were not judged. The synthetic data was
written with the kinds of conversation and question in these two
benchmarks in mind, so the gains say nothing about data of another
kind. The LongMemEval-S figures are from one schedule. The model is
English only. For these reasons, and because of the cost above, the
static embedder remains the default; the fine-tuned one suits a store
that is written slowly and searched often, or a machine with a GPU.
The weights are published under the MIT licence as
\texttt{myelinagent/myelin-embed-small-en-v1}, with the data and
training script in the repository.

\section{Limitations}

\begin{itemize}
\item \textbf{Self-run benchmarks.} All scores are ours, on two public
  benchmarks, with one to four runs per setting. Other systems
  were run only on the static arm of the over-time protocol, on
  LoCoMo; there, with the same static embedder, Mem0
  retrieved as much of the evidence in the top 5 as \mbox{Myelin-AI}
  without feedback (47.2\% against 46.4\%) and more in the top 10 and
  20, at one model call per turn written. Details are in the companion
  benchmark paper.
\item \textbf{Reader-dependent accuracy.} The answering model and
  prompt account for more of the benchmark score than the memory does.
  The memory's own contribution is better read from the retrieval
  figures.
\item \textbf{Forgetting is tuned on one dataset} and fails on
  another, as described above.
\item \textbf{Linear vector search.} Hybrid retrieval is exact and
  scans every vector.
\item \textbf{Lexical dependence.} Excerpts and the lexical ranking
  rely on shared words; a question phrased with none of the evidence's
  vocabulary depends on the embedder alone.
\item \textbf{Single-process store.} One SQLite file per deployment;
  concurrency and multi-tenant operation at scale are untested.
\item \textbf{Unevaluated components.} The salience gate and
  consolidation routine are not measured here.
\end{itemize}

\section{Conclusion}

\mbox{Myelin-AI} shows that a memory layer can be small, model-free on its write
path and auditable, and still sit in the range of published results on
the standard benchmarks. Its more distinctive parts (feedback from
use, relationships, and reversible forgetting from one grade) are
invisible to those benchmarks and measurable over time: feedback
reliably helps, and forgetting is safe only where use arrives before
decay does. We release the library, the harnesses and every result
file so that both findings can be checked.

\paragraph{Availability.} \texttt{pip install myelin-ai}; source,
benchmark harnesses and result files are in the project repository
under the MIT licence. LoCoMo \citep{maharana2024locomo} is distributed
by its authors under CC BY-NC 4.0 and LongMemEval
\citep{wu2025longmemeval} under the MIT licence; neither is
redistributed with this work, and both were used unmodified for
evaluation.

\bibliographystyle{plainnat}
\bibliography{refs}
\end{document}
